\begin{papercard}{obrien2014}
\cf{Stages} \stagetag{R} (what to report about a qualitative study), \stagetag{C} (reporting the analysis, unit of analysis, instruments and coding team), \stagetag{I} (how multiple coders and discrepancies are described), \stagetag{A} (sampling of documents). The PDF in the corpus is \emph{only the Supplemental Digital Appendices} to O'Brien, Harris, Beckman, Reed and Cook, ``Standards for reporting qualitative research: a synthesis of recommendations'' (\emph{Academic Medicine}, 2014): Appendix~1 lists the sources used, Appendix~2 explains and illustrates each of the 21 SRQR items, followed by a reference list (45 PDF pages). The main article (methods of the synthesis, Table~2 mapping the items to the original sources) is \emph{not} in the file, so the card does not describe how the standards were derived beyond what the appendix says.

\cf{Appendix 1: sources} The full list holds 40 sources used to develop the standards; of these, 25 are the most complete sources listed in the article's Table~2 and 15 are additional \pc{obrien2014}{1}{This Appendix includes the 25 sources listed in Table 2 as well as the additional 15}. The list mixes checklists and appraisal tools (CASP qualitative checklist, the AERA standards for empirical social science research, Elliott et al.\ guidelines for psychology, Giacomini and Cook's users' guides, Malterud, Mays and Pope, Popay et al.) with the COREQ checklist for interviews and focus groups \pc{obrien2014}{6}{Consolidated criteria for reporting qualitative research}. It shows that many overlapping standards existed and were synthesised rather than newly invented.

\cf{How the appendix says the items should be used} Appendix~2 gives explanations and examples for the 21 SRQR items \pc{obrien2014}{7}{Explanations and Examples of the 21 Items from the Standards for Reporting Qualitative}. Authors may not be able to report every element, so they should ``prioritize elements that are most relevant to the given study, findings, context, and readership''; journal conventions may change the order \pc{obrien2014}{7}{prioritize elements that are most relevant to the given study, findings, context, and readership}. Examples were drawn from 2011--2013 papers in four medical-education journals \pc{obrien2014}{7}{We identified recent (2011-2013) examples from journals that frequently publish qualitative}.

\cf{The 21 items} \emph{Title and abstract}: 1 title (state that the study is qualitative or name the approach or data-collection method); 2 abstract (background, purpose, methods, findings, implications; ``findings'' may need translating into ``results'' for structured abstracts). \emph{Introduction}: 3 problem formulation, with the gap and the scope of what will and will not be included \pc{obrien2014}{11}{It should provide an overview of what is known about the problem, highlight gaps in current knowledge}; 4 purpose or research question (qualitative studies often ask ``how'' and ``why''). \emph{Methods}: 5 qualitative approach and research paradigm; 6 researcher characteristics and reflexivity; 7 context; 8 sampling strategy; 9 ethical issues; 10 data collection methods; 11 data collection instruments and technologies; 12 units of study; 13 data processing; 14 data analysis; 15 techniques to enhance trustworthiness. \emph{Results}: 16 synthesis and interpretation; 17 links to empirical data. \emph{Discussion}: 18 integration with prior work, implications, transferability and contribution; 19 limitations. \emph{Other}: 20 conflicts of interest; 21 funding.

\cf{Items 5 and 6: paradigm and reflexivity} Authors should name the paradigm because it tells readers ``whether the researcher assumes that there is a single, objective reality'' (positivist or post-positivist) or multiple subjective realities (constructivist or interpretivist) and has consequences for design and rigour \pc{obrien2014}{13}{whether the researcher assumes that there is a single, objective reality}; the approach (content analysis, general inductive approach, grounded theory, ethnography) must be named as well, with a rationale and references. Reflexivity is the ``intentional, systematic consideration of the potential or actual effects of the researcher(s) on all aspects of the study'' \pc{obrien2014}{15}{Reflexivity refers to intentional, systematic consideration of the potential or actual effects of the}; in post-positivist work researcher characteristics are seen as possible bias, in interpretivist work as accepted context; authors should state background, training, stance, assumptions, and relationship to participants. They also note that ``there is no expectation that the study could be precisely replicated'' and that such characteristics should not go in the limitations section \pc{obrien2014}{16}{There is no expectation that the study could be precisely replicated}.

\cf{Items 7--9: context, sampling, ethics} Context: the setting and the cultural, political, historical, economic or other external factors that situate the findings \pc{obrien2014}{17}{describe the setting/site(s) in which the study was conducted}. Sampling: describe \emph{how} documents, events or participants were selected instead of only labelling the strategy ``purposive'' or ``snowball'' \pc{obrien2014}{18}{We recommend that authors describe the sampling strategy rather than simply labeling the strategy}; give the criteria used to stop sampling and, if ``saturation'' is invoked, define it \pc{obrien2014}{19}{then the specific criteria used to define saturation should be described}; sample size need not be fixed beforehand but must be explained afterwards \pc{obrien2014}{19}{nonetheless describe how they established the final sample size}. Ethics: approval from a review board, consent, privacy and data security, anonymisation of excerpts, and any incentives \pc{obrien2014}{21}{If researchers provided compensation or offered incentives to facilitate participation, this should be reported}.

\cf{Items 10--13: data collection and units} Data collection is iterative: analysis often begins before collection ends, so authors ``must clearly describe the iterative process of data collection and analysis'' and any changes of procedure in response to emerging findings \pc{obrien2014}{22}{authors must clearly describe the iterative process of data collection and analysis}; the study period (start and end dates of collection and analysis) is to be reported \pc{obrien2014}{22}{The study period (start and end dates for data collection and analysis) should be identified}. Instruments include interview guides, observation protocols and ``data extraction or coding protocols for selection and analysis of documents, photographs, videos, or other artifacts''; authors should describe how they were developed and changed and ideally give access \pc{obrien2014}{24}{It is often helpful for authors to provide access to the data collection instrument(s) or a detailed description of them}. Item 12 separates the \emph{planned} sample (item 8) from the \emph{actual} units: ``Item 12 focuses on description of the actual participants, documents or events included'' and how they differ from the target \pc{obrien2014}{26}{By contrast, Item 12 focuses on description of the actual}; for documents report source, intended audience, date and type \pc{obrien2014}{26}{For documents, this might include the source, intended audience, date, or type of document}. Item 13: preparation, management, verification of integrity, and when anonymisation was applied \pc{obrien2014}{28}{Authors should describe procedures used to maintain data security and protect the privacy of}.

\cf{Item 14: data analysis (most relevant to coding)} Authors should be transparent about the process of analysis so readers can follow the logic from question to finding, cite the guiding approach if one exists, and justify deviations. They ``should specify the unit of analysis'' \pc{obrien2014}{29}{Authors should specify the unit of analysis}, which need not equal the unit of sampling \pc{obrien2014}{29}{the unit of analysis is not necessarily the same as the unit of sampling}. Discrepant cases that do not fit the identified concepts must be handled and described: ``the authors must describe how these discrepancies were handled during the analysis'' \pc{obrien2014}{30}{the authors must describe how these discrepancies were handled during the analysis}. Theoretical frameworks that shaped the categories (``sensitizing concepts'') must be named; the team members involved and their perspectives given; and software use specified: ``Simply stating that software was used is insufficient'' \pc{obrien2014}{31}{Simply stating that software was used is insufficient}. Examples include independent coding followed by discussion of open codes, a consensus-building approach, an analytic matrix audited by two researchers outside the coding, and a ``deviant case'' check; one example reaches ``full agreement'' through discussion, with no agreement statistic.

\cf{Item 15: trustworthiness} Techniques named: member checking, triangulation of sources, methods or researchers, an explicit audit trail, prolonged immersion \pc{obrien2014}{33}{member checking; triangulation of data sources, methods, and/or researchers; creation of an explicit audit trail}. An audit trail is ``careful documentation of all decisions made throughout the study'' so that an external researcher can review every step \pc{obrien2014}{33}{An audit trail involves careful documentation of all decisions made throughout the study}; authors must explain why the chosen techniques suit the study.

\cf{Items 16--17: findings and evidence} In qualitative work the line between results and discussion ``tends to blur'' \pc{obrien2014}{35}{In qualitative research the distinction between results and discussion tends to blur}. Frequency counts ``play a limited role'' and need not be reported unless they matter for interpretation \pc{obrien2014}{36}{or codes) play a limited role in qualitative research, and need not be reported unless they play a meaningful role}. Findings need evidence: quotes, field notes, excerpts or photographs, de-identified, with an appendix or repository if space is short \pc{obrien2014}{37}{an appendix, supplemental material, or web-based repository could be used to provide access to additional data}.

\cf{Items 18--21: discussion and other} Start with a short summary of main findings, connect them to prior work and state transferability \pc{obrien2014}{39}{The authors should begin the Discussion with a short summary of the main findings}. Limitations: techniques for trustworthiness belong in the methods, whereas the discussion identifies gaps in those efforts and consequences for transferability \pc{obrien2014}{40}{In the Discussion, authors should identify problems or gaps in their efforts to ensure trustworthiness}. Conflicts of interest and how they were managed \pc{obrien2014}{41}{Authors should also explain how these conflicts were managed in the conduct of the study}; funding and funders' roles \pc{obrien2014}{42}{The authors should describe any sources of funding and other support for the study}.

\cf{Limits (card)} Only the appendices are in the corpus: no derivation, no item-to-source mapping, no evaluation of whether SRQR improves reporting. Examples come from medical-education interview and focus-group studies, so document-based coding is only covered through brief mentions. The appendix names no inter-coder reliability statistic (the only agreement mention is ``full agreement was reached'' by discussion) and treats counts as marginal, which does not match a codebook study that reports $\kappa$ or $\alpha$. The reference list numbering is local to the appendix, as the authors state \pc{obrien2014}{43}{The numbering of the references in this Appendix does not match that of the references in the}.

\cf{Use for this thesis} (1) A reporting scaffold for the qualitative parts of the thesis (codebook development, coding of policies and walkthrough material): item 5 forces a stated paradigm and approach (for example codebook-based content analysis with a post-positivist stance versus an interpretive one), to be settled in the Decision Log; item 6 asks for the coder's background and stance (one human coder plus, if used, an LLM). (2) Items 8 and 12 give the vocabulary to separate the \emph{planned} app and policy sample from the \emph{actual} one, and the document attributes to record: source, intended audience, date (version of the policy), type. (3) Item 11 treats coding protocols as instruments: the codebook, its changes during piloting and its availability belong in the report. (4) Item 14: name the unit of analysis (policy segment, data flow, app, screen), the handling of disagreements and ``deviant'' cases, the software, and the theoretical frameworks behind the categories (Contextual Integrity, Citron and Solove). For an LLM the equivalent of ``software was used'' is model, version, prompt, settings and date. (5) Item 15: audit trail, researcher triangulation (second coder), and archived snapshots of policies and screens as the evidence for item 17. (6) Complement with a quantitative reliability report (Krippendorff, Lombard, Artstein and Poesio, McHugh) because SRQR gives no statistic.
\end{papercard}
